\documentclass[11pt,a4paper]{article}
\usepackage[T1]{fontenc}
\usepackage[utf8]{inputenc}
\usepackage[brazilian]{babel}
\usepackage{lmodern}
\usepackage[margin=2.1cm]{geometry}
\usepackage{enumitem}
\usepackage{xcolor}
\usepackage{amssymb}
\usepackage{fancyhdr}

\pagestyle{fancy}
\fancyhf{}
\lhead{\small TCC, César Hoffmann}
\rhead{\small UNISINOS \textbullet{} Perfil A}
\cfoot{\small\thepage}
\renewcommand{\headrulewidth}{0pt}
\setlength{\parindent}{0pt}
\setlength{\parskip}{0.4em}
\setlist[enumerate]{label=\textbf{\arabic*.},leftmargin=*,itemsep=0.55em,topsep=0.2em}
\newcommand{\aprofundar}[1]{\par{\small\color{gray}\textit{Se couber: #1}}}
\newcommand{\marcar}{$\square$\ }

\begin{document}
\begin{center}
  {\Large\bfseries Entrevista com gestor, perfil A}\\[0.3em]
  {\small Estabelecimento anterior a 2020 que adotou autoatendimento digital}\\[0.3em]
  Cardápios digitais em bares e restaurantes de Porto Alegre
\end{center}

\textbf{Casa:} \rule{7cm}{0.4pt}\quad
\textbf{Data/local:} \rule{4cm}{0.4pt}

\textbf{Pessoa e função:} \rule{9cm}{0.4pt}

\textbf{Antes de gravar:} \marcar levar o TCLE em duas vias e colher a assinatura; \marcar confirmar a gravação e a opção de identificação da casa; \marcar pedir autorização específica para fotos ou documentos. Não divulgar o nome pessoal. Perguntas podem ser puladas. Não solicitar contratos, dados de clientes nem senhas.

\textbf{Tempo:} cerca de 30 a 45 minutos; reservar até 60 se houver disponibilidade. Perguntar primeiro, depois pedir um exemplo real: quando ocorreu, o que aconteceu e como foi resolvido. Distinguir o que a pessoa viu do que o fornecedor informou.

\section*{1. A casa e a mudança}
\begin{enumerate}
  \item Quando a casa abriu e qual é o seu papel na operação?
    \aprofundar{Número aproximado de mesas, equipe e perfil dos clientes, se a pessoa souber.}
  \item Como o cliente consultava o cardápio e fazia o pedido antes de 2020? E como faz hoje?
    \aprofundar{Separar PDF, página para consulta e pedido feito pelo próprio cliente; salão e delivery.}
  \item Em que momento passaram a oferecer o pedido digital ao cliente? O que levou a essa decisão?
    \aprofundar{Quem decidiu? Foi uma mudança planejada ou feita durante a pandemia?}
  \item O que mudou na rotina quando passaram a usar essa solução?
    \aprofundar{Pedir um episódio concreto, positivo ou negativo, sem supor ganho de tempo ou redução de equipe.}
\end{enumerate}

\section*{2. O sistema na prática}
\begin{enumerate}[resume]
  \item Como o cliente faz o pedido digital, do acesso ao cardápio até a confirmação?
    \aprofundar{QR Code, tablet, totem ou outro; como ele sabe que o pedido chegou.}
  \item Qual solução vocês contrataram ou desenvolveram? Por que escolheram essa?
    \aprofundar{Houve outra opção ou troca de fornecedor? O que sabem por experiência própria?}
  \item Depois de confirmado, como o pedido chega ao caixa e à cozinha?
    \aprofundar{Existe integração que vocês usam de fato ou alguém precisa repetir a informação?}
  \item Quem atualiza preços e itens esgotados? Quanto tempo demora para a mudança aparecer ao cliente?
    \aprofundar{Lembram de alguma vez em que o cardápio mostrava um item indisponível?}
  \item Já faltou internet ou o sistema parou em dia de atendimento? O que vocês fizeram?
    \aprofundar{Cardápio físico, anotação manual, outro link, suporte; pedir um caso real.}
\end{enumerate}

\section*{3. Qualidade e necessidades}
\begin{enumerate}[resume]
  \item O que costuma dar mais trabalho ou atrapalhar o uso do sistema?
    \aprofundar{Carregamento, atualização, estabilidade, cadastro, leitura no celular, suporte.}
  \item O que é indispensável em um cardápio digital para a casa funcionar bem?
    \aprofundar{Pedir até três prioridades; que problema cada uma resolve?}
  \item O que os clientes comentam sobre o pedido digital e sobre a opção de cardápio físico?
    \aprofundar{Distinguir comentário ouvido de impressão da equipe; pedir exemplo sem identificar cliente.}
  \item O sistema pede dados pessoais para consultar ou pedir? Para quê?
    \aprofundar{Distinguir pedido no salão, delivery e fidelidade; não solicitar dados pessoais.}
\end{enumerate}

\section*{4. Conferência de requisitos e encerramento}
\begin{enumerate}[resume]
  \item Para vocês, o que seria mais importante melhorar: atualização dos itens, funcionamento sem internet, rapidez ou leitura da tela? Algum desses pontos não é um problema aqui?
    \aprofundar{Perguntar a prioridade e o motivo antes de mostrar soluções propostas no TCC.}
  \item Existe algo que um cardápio digital não deve fazer, mesmo que pareça prático para o cliente?
  \item Tem alguma situação importante que eu não perguntei? Posso pedir depois um esclarecimento sobre um trecho da entrevista?
\end{enumerate}

\textbf{Depois:} anotar pontos contraditórios, exemplos e o que foi visto diretamente; registrar se fotos foram autorizadas. Guardar o áudio e o TCLE fora do repositório do trabalho.

\vspace{0.6em}
\textbf{Anotações:}\par
\newcount\linhas \linhas=0
\loop\ifnum\linhas<7 \rule{\linewidth}{0.25pt}\par\vspace{0.55em}\advance\linhas by 1 \repeat
\end{document}
